\documentclass[a4paper,11pt]{article}
\usepackage[a4paper,margin=2cm]{geometry}
\usepackage[T1]{fontenc}
\usepackage[utf8]{inputenc}
\usepackage[ngerman,english]{babel}
\usepackage{lmodern}
\usepackage{microtype}
\usepackage{xcolor}
\usepackage{parskip}
\usepackage{enumitem}
\usepackage[most]{tcolorbox}
\usepackage{titlesec}
\usepackage[hidelinks]{hyperref}
\definecolor{HeaderBlueDark}{RGB}{70,110,170}
\definecolor{RowBlueLight}{RGB}{235,243,255}
\definecolor{RowBlueDark}{RGB}{220,233,250}
\definecolor{SoftGray}{RGB}{90,90,90}
\setlength{\parindent}{0pt}
\setlist[itemize]{leftmargin=1.4em,itemsep=0.35em,topsep=0.35em}
\linespread{1.06}
\titleformat{\section}[block]{\Large\bfseries\color{HeaderBlueDark}}{}{0pt}{}
\titlespacing{\section}{0pt}{1.3em}{0.8em}
\titleformat{\subsection}[block]{\large\bfseries\color{HeaderBlueDark}}{}{0pt}{}
\titlespacing{\subsection}{0pt}{1.0em}{0.6em}
\newtcolorbox{exbox}{enhanced,breakable,colback=RowBlueLight,colframe=RowBlueDark,boxrule=0.4pt,arc=1.5mm,left=1.2mm,right=1.2mm,top=1mm,bottom=1mm,title={Beispiele \textnormal{(Examples)}},fonttitle=\bfseries,colbacktitle=RowBlueDark,coltitle=black}
\NewDocumentEnvironment{grammarentry}{mm+b}{%
\begin{tcolorbox}[enhanced,breakable,colback=white,colframe=HeaderBlueDark,boxrule=0.6pt,arc=1.5mm,left=1.5mm,right=1.5mm,top=1mm,bottom=1mm,colbacktitle=HeaderBlueDark,coltitle=white,fonttitle=\bfseries,title={#1\hfill\normalfont\footnotesize #2}]
#3
\end{tcolorbox}}{}
\newcommand{\GermanText}[1]{\textbf{Deutsch: }#1\par}
\newcommand{\EnglishText}[1]{{\small\color{SoftGray}\textbf{English: }#1\par}}
\newcommand{\ExampleItem}[2]{\item \textbf{#1}\\{\small\color{SoftGray}#2}}


\title{W-Relativpronomen: \glqq wessen\grqq{}}

\date{}

\begin{document}

\maketitle

\tableofcontents

\newpage

\section{Grundbedeutung und grammatische Rolle}

\begin{grammarentry}{wessen}{W-Relativpronomen, Genitiv}

\GermanText{\textbf{wessen} kann einen \textbf{freien Relativsatz} einleiten. Es bedeutet ungefähr \glqq whose; whoever's\grqq{}.}

\EnglishText{\textbf{wessen} can introduce a \textbf{free relative clause}. It means approximately \glqq whose; whoever's\grqq{}.}

\GermanText{Im Relativsatz hat \textbf{wessen} die Funktion \textbf{Genitivattribut / Besitz}.}

\EnglishText{Inside the relative clause, \textbf{wessen} functions as \textbf{Genitivattribut / Besitz}.}

\GermanText{\textbf{wessen} ist die Genitivform der Reihe \textbf{wer - wen - wem - wessen}.}

\EnglishText{\textbf{wessen} is the genitiv form in the series \textbf{wer - wen - wem - wessen}.}

\GermanText{Es steht normalerweise direkt vor einem Nomen und drückt Zugehörigkeit oder Besitz aus. Die Form \textbf{wessen} bleibt dabei unverändert.}

\EnglishText{It normally stands directly before a noun and expresses belonging or possession. The form \textbf{wessen} remains unchanged.}

\end{grammarentry}

\section{Freier Relativsatz}

\begin{grammarentry}{Kein ausdrückliches Bezugsnomen}{Bedeutung bereits enthalten}

\GermanText{Vor \textbf{wessen} steht normalerweise kein konkretes Bezugsnomen. Die Bedeutung ist ungefähr \glqq die Person, die ...\grqq{} oder \glqq jeder, der ...\grqq{}.}

\EnglishText{There is normally no concrete antecedent noun before \textbf{wessen}. The meaning is approximately \glqq the person who ...\grqq{} or \glqq whoever ...\grqq{}.}

\GermanText{Wenn bereits ein Nomen vorhanden ist, nimmt man normalerweise \textbf{der/die/das} und die passenden deklinierten Formen.}

\EnglishText{If a noun is already present, German normally uses \textbf{der/die/das} and the appropriate declined forms.}

\end{grammarentry}

\section{Beispiele}

\begin{exbox}

\begin{itemize}

\ExampleItem{Wessen Auto hier steht, der soll es wegfahren.}{Whoever's car is parked here should move it.}

\ExampleItem{Wessen Name auf der Liste steht, darf eintreten.}{Whoever's name is on the list may enter.}

\ExampleItem{Wessen Hilfe du brauchst, den solltest du fragen.}{You should ask the person whose help you need.}

\end{itemize}

\end{exbox}

\section{Typischer Fehler}

\begin{grammarentry}{Nicht mit einem normalen Bezugsnomen}{Korrektur}

\GermanText{Falsch: \textbf{Wem Auto hier steht, der soll kommen.}}

\EnglishText{Incorrect: Wem Auto hier steht, der soll kommen.}

\GermanText{Richtig: \textbf{Wessen Auto hier steht, der soll kommen.}}

\EnglishText{Correct: Wessen Auto hier steht, der soll kommen.}

\end{grammarentry}

\section{Wortstellung und Komma}

\begin{grammarentry}{Nebensatzregel}{Verb am Ende}

\GermanText{Der freie Relativsatz ist ein Nebensatz. Das konjugierte Verb steht normalerweise am Ende. Der Relativsatz wird durch Kommas abgetrennt.}

\EnglishText{A free relative clause is a subordinate clause. The conjugated verb normally stands at the end. The relative clause is separated by commas.}

\end{grammarentry}

\section{Merksatz}

\begin{grammarentry}{wessen}{kurz und wichtig}

\GermanText{\textbf{wessen} = Genitivform eines W-Relativpronomens für Personen in freien Relativsätzen.}

\EnglishText{\textbf{wessen} = the genitiv form of a W-relative pronoun for people in free relative clauses.}

\end{grammarentry}

\end{document}